\chapter{Introduction}

\section{Background}

Tissue engineering is a relatively new field within the multidisciplinary realm of biomedical engineering. It combines the fields of biology, medicine, material science, and engineering into a single discipline in which three-dimensional structures are fabricated with the intention of restoring, maintaining, or even improving the function of tissue. The art of engineering tissue has many applications within medicine, cosmetics, and cellular research, and there are many methods of fabricating these tissue scaffolds, with one common approach being 3D printing.

The ability to create structures that replicate tissue functions, whether seeded with cells or not, has countless applications within various fields such as medicine, cosmetics, and biological research. It creates the potential to manufacture substitute tissues (or even organs in more advanced applications) in regenerative medicine for transplantation to mimic the function of damaged or missing tissues in patients. It is often used in drug development and testing. Tissue engineering is also used in biological research, where 3D environments are created to simulate and study the growth and behaviour of cellular systems such as cancer cells and bacteria.

A popular class of materials used to fabricate substitute tissue structures are hydrogels, which can be processed in various ways, with 3D printing being a particularly popular method. Hydrogels are a broad class of materials that are known for their tunable properties, porous nature, and ability to retain large amounts of water, which allows many hydrogels to support cell growth. However, due to the complex rheological properties of different hydrogels and the varying methods by which these materials are hardened, it remains a challenge to 3D print complex structures in a way that is accurate, repeatable, and cost-effective. This project focuses on the design, construction, and testing of a hydrogel extrusion system capable of creating 3D scaffolds to overcome these challenges.

\section{Objectives}

The primary objectives of this project are:
\begin{enumerate}
    \item To design and assemble a hydrogel extruder capable of controlled and stable extrusion.
    \item To design and construct a three-axis gantry to position the extruder for accurate 3D printing.
    \item To design and implement an embedded system (both hardware and software) capable of real-time control of system actuators, sensing, and closed-loop logic.
    \item To test and validate the mechanical and embedded systems, and to design an experimental method for determining optimal printing parameters for different bioinks.
\end{enumerate}

The scope of the project is limited to developing and demonstrating a working prototype. While the system is not intended to achieve the full performance and quality of commercial bioprinters, it provides a validated platform for future upgrades and serves as a tool for characterising bioinks.

\section{Motivation}

The motivation for this project lies in the growing importance of tissue engineering and the need for accessible tools to support research in this field. Most commercial bioprinters remain prohibitively expensive and proprietary, limiting their availability to smaller research groups. Developing a low-cost prototype capable of producing controlled, stable, and repeatable prints allows research in tissue engineering to be conducted by a wider range of researchers. This project therefore contributes to making tissue engineering more affordable and provides a means of characterising different gels, offering insight into their behaviour in both tissue engineering and other biomedical applications.